\paragraph{\texorpdfstring{$S_3$}{S3} 预解式塔的 endoscopic 分解、动机直和与端同态刚性}\label{par:xi-terminal-zm-s3-endoscopic-motive-polarization-endomorphism-rigidity}
记
\[
R:=\overline{\mathcal R},\qquad
\mathscr R(z,y)=z^3+(1+2y)z^2-(1+4y+4y^2)z-(1+5y+13y^2+8y^3).
\]
由结论 \ref{con:xi-terminal-zm-resolvent-cubic-plane-cubic-geometry}，\(R\subset\mathbf A^2_{z,y}\) 是定义在 \(\QQ\) 上的非奇异平面三次曲线，因而为椭圆曲线；其到 \(y\)-直线的投影为次数 \(3\) 覆盖，分歧由 \(\Disc_z(\mathscr R)\) 控制。
同时，\(C_6\) 可由仿射方程组
\[
w^2=-y(y-1)P_{\mathrm{LY}}(y),\qquad \mathscr R(z,y)=0
\]
经规范化得到，其中
\[
P_{\mathrm{LY}}(y)=256y^3+411y^2+165y+32.
\]
由结论 \ref{con:xi-terminal-zm-discriminant-ridge-etale-c3-explicit-model}，存在无分歧循环三次覆盖
\[
\pi:C_6\to C,\qquad \mathrm{Gal}(C_6/C)\cong C_3,
\]
且复合映射 \(C_6\to\PP^1_y\) 为 \(S_3\)-Galois 覆盖。

\begin{theorem}[\texorpdfstring{$S_3$}{S3} 预解式塔的 endoscopic 同源分解与二维标准因子识别]\label{thm:xi-terminal-zm-s3-endoscopic-homology-a2-identification}
在上述设定下，令
\[
P:=\operatorname{Prym}(C_6/C):=\ker\!\bigl(\operatorname{Nm}_{\pi}:J(C_6)\to J(C)\bigr)^0.
\]
则存在定义在 \(\QQ\) 上的同源分解
\[
J(C_6)\sim_{\QQ}J(C)\times J(R)^2,
\qquad
P\sim_{\QQ}J(R)^2.
\]
特别地，先前由 \(S_4\) 或 \(S_3\) 表示分解得到的二维标准因子（记作 \(A_{2\text{-}\mathrm{dim}}\) 或 \(A_2\)）在 \(\QQ\)-同源意义下可精确识别为
\[
A_{2\text{-}\mathrm{dim}}\sim_{\QQ}J(R)^2.
\]
\leanverified{paper\_xi\_terminal\_zm\_s3\_endoscopic\_homology\_a2\_identification}
\end{theorem}
\begin{proof}
由
\(
\QQ[S_3]\cong \QQ\oplus\QQ\oplus M_2(\QQ)
\)
与图对应作用，\(J(C_6)\) 在 \(\QQ\)-同源意义下分解为平凡块、符号块与标准块。
因 \(C_6/S_3\simeq\PP^1\)（亏格 \(0\)），平凡块在 Jacobian 层面消失。
又 \(C_6/A_3\simeq C\)，故符号块对应 \(J(C)\)。

取任一换位 \(\tau\in S_3\)，由
\(
C_6/\langle\tau\rangle\simeq R
\)
得 \(J(R)\) 作为 \(J(C_6)\) 的同源子簇出现。
同时换位在符号表示上作用为 \(-1\)，故该子簇不落入符号块，只能落入标准块 \(P\)。
因此 \(P\) 含有一个同源于 \(J(R)\) 的椭圆因子。

另一方面，标准块上有
\(
M_2(\QQ)\hookrightarrow \End^0_{\QQ}(P)
\)。
对二维阿贝尔簇，若有 \(M_2(\QQ)\) 作用，则秩 \(1\) 幂等元给出椭圆投影，故该簇必同源于某椭圆曲线平方。
结合上一步中已出现的椭圆因子 \(J(R)\)，可得
\(
P\sim_{\QQ}J(R)^2
\)。
与符号块合并即得
\(
J(C_6)\sim_{\QQ}J(C)\times J(R)^2
\)。
\end{proof}

\begin{corollary}[有限域点数恒等式与局部 \texorpdfstring{$\zeta$}{zeta} 因子平方分解]\label{cor:xi-terminal-zm-s3-endoscopic-pointcount-localzeta-square}
\leanverified{paper\_xi\_terminal\_zm\_s3\_endoscopic\_pointcount\_localzeta\_square}
在定理 \ref{thm:xi-terminal-zm-s3-endoscopic-homology-a2-identification} 的设定下，取同时使 \(C,C_6,R\) 良好约化且 \(p\nmid 6\) 的素数 \(p\)。
则对任意 \(n\ge 1\) 成立
\[
\#C_6(\FF_{p^n})
=\#C(\FF_{p^n})+2\,\#R(\FF_{p^n})-2(p^n+1).
\]
等价地，若
\[
Z(X/\FF_p,T)=\frac{P_{X,p}(T)}{(1-T)(1-pT)},
\]
则
\[
P_{C_6,p}(T)=P_{C,p}(T)\,P_{R,p}(T)^2.
\]
\end{corollary}
\begin{proof}
由定理 \ref{thm:xi-terminal-zm-s3-endoscopic-homology-a2-identification}，
\[
H^1_{\acute et}(C_{6,\overline{\FF}_p},\QQ_\ell)
\cong
H^1_{\acute et}(C_{\overline{\FF}_p},\QQ_\ell)
\oplus
H^1_{\acute et}(R_{\overline{\FF}_p},\QQ_\ell)^{\oplus2}
\]
在好素数处与 Frobenius 作用相容。
特征多项式乘法给出局部 \(P\)-因子分解；对 Frobenius 幂次取迹并代回 Weil 点数公式即得点数恒等式。
\end{proof}

\begin{theorem}[Chow 动机范畴中的 \texorpdfstring{$h^1$}{h1} 正交分解]\label{thm:xi-terminal-zm-s3-endoscopic-chow-motive-h1-splitting}
在带 \(\QQ\)-系数的 Chow 动机伪阿贝尔闭包范畴中，存在同构
\[
h^1(C_6)\simeq h^1(C)\oplus h^1(R)^{\oplus2}.
\]
该分解可由 \(S_3\) 作用诱导的群代数幂等元对应显式实现；其任意 \(\ell\)-进实现即为推论
\ref{cor:xi-terminal-zm-s3-endoscopic-pointcount-localzeta-square} 中的 \(H^1_{\acute et}\) 直和分解。
\end{theorem}
\begin{proof}
\(S_3\) 对 \(C_6\) 的作用给出
\[
\QQ[S_3]\hookrightarrow \mathrm{Corr}^0(C_6,C_6)\otimes\QQ.
\]
取中心幂等元得到 \(h^1(C_6)\) 的正交分块。
平凡块对应商 \(C_6/S_3\simeq\PP^1\)，故在 \(h^1\) 上为零；
符号块通过商映射 \(C_6\to C_6/A_3\simeq C\) 的 pull--push 对应与 \(h^1(C)\) 同构。

标准块对应 \(M_2(\QQ)\)。
在该块内取两条互补秩 \(1\) 幂等元，可分别得到两块同构的动机像。
由换位商 \(C_6/\langle\tau\rangle\simeq R\)（任一换位 \(\tau\)）可识别每一像为 \(h^1(R)\)，故标准块同构于 \(h^1(R)^{\oplus2}\)。
综上得到所述分解。
\end{proof}

\leanverified{paper\_xi\_terminal\_zm\_s3\_endoscopic\_prym\_a2\_coxeter}
\begin{proposition}[Prym 极化的 \texorpdfstring{$A_2$}{A2} 格矩阵与 deck 变换的 Coxeter 表示]\label{prop:xi-terminal-zm-s3-endoscopic-prym-a2-coxeter}
设 \(P=\operatorname{Prym}(C_6/C)\)。
则存在定义在 \(\QQ\) 上的同源
\[
\varphi:P\longrightarrow J(R)^2
\]
使得 \(P\) 上 Prym 极化 \(\lambda_P\) 在 \(J(R)^2\) 上对应的 Riemann 形式在适当整基下可取为
\[
Q_{A_2}=
\begin{pmatrix}
2 & -1\\
-1 & 2
\end{pmatrix}.
\]
因此 \(\det(Q_{A_2})=3\)，极化型为 \((1,3)\)。
并且可取三次 deck 变换在该基下由
\[
U=
\begin{pmatrix}
0 & -1\\
1 & -1
\end{pmatrix}\in\mathrm{SL}_2(\ZZ)
\]
表示，满足
\[
\operatorname{tr}(U)=-1,\qquad \chi_U(T)=T^2+T+1.
\]
\end{proposition}
\begin{proof}
矩阵 \(Q_{A_2}\) 与极化型 \((1,3)\) 已由定理
\ref{thm:xi-terminal-zm-prym-polarization-a2-rigidity} 给出。
另一方面，命题 \ref{prop:xi-s3-closure-standard-isotypic-explicit-r3-realization} 中三循环 \((123)\) 在标准块上的矩阵恰可取为
\(
\begin{psmallmatrix}0&-1\\1&-1\end{psmallmatrix}
\)，其特征多项式即 \(T^2+T+1\)。
该矩阵保持 \(Q_{A_2}\)，对应 \(A_2\) 根格的 Coxeter 元表示。
\end{proof}

\begin{theorem}[\texorpdfstring{$J(C_6)$}{J(C6)} 的绝对端同态代数与 Picard 数二值刚性]\label{thm:xi-terminal-zm-s3-endoscopic-absolute-endomorphism-picard-dichotomy}
设 \(J(C)=\operatorname{Jac}(C)\) 满足结论
\ref{con:xi-terminal-zm-discriminant-ridge-jacobian-absolute-simplicity-endomorphism}，
即 \(J(C)\) 在 \(\overline{\QQ}\) 上绝对单且
\(
\End^0_{\overline{\QQ}}(J(C))\cong\QQ
\)。
则
\[
\End^0_{\overline{\QQ}}(J(C_6))
\cong
\QQ\times M_2\!\bigl(\End^0_{\overline{\QQ}}(J(R))\bigr).
\]
从而存在严格二分：
\begin{enumerate}
  \item 若 \(J(R)\) 非 CM，则
  \[
  \End^0_{\overline{\QQ}}(J(C_6))\cong \QQ\times M_2(\QQ).
  \]
  \item 若 \(J(R)\) 为 CM 椭圆曲线（CM 域 \(K\)），则
  \[
  \End^0_{\overline{\QQ}}(J(C_6))\cong \QQ\times M_2(K).
  \]
\end{enumerate}
并且几何 N\'eron--Severi 秩满足
\[
\rho\!\bigl(J(C_6)_{\overline{\QQ}}\bigr)=
\begin{cases}
4,& J(R)\ \text{非 CM},\\
5,& J(R)\ \text{CM}.
\end{cases}
\]
\end{theorem}
\begin{proof}
由定理 \ref{thm:xi-terminal-zm-s3-endoscopic-homology-a2-identification}，
\[
J(C_6)\sim_{\overline{\QQ}}J(C)\times J(R)^2.
\]
又因 \(J(C)\) 绝对单且维数 \(2\)，不存在非零同态 \(J(C)\to J(R)\)；故
\(
\Hom_{\overline{\QQ}}(J(C),J(R))=0
\)。
因此端同态代数按直积分解：
\[
\End^0_{\overline{\QQ}}(J(C_6))
\cong
\End^0_{\overline{\QQ}}(J(C))
\times
\End^0_{\overline{\QQ}}(J(R)^2)
\cong
\QQ\times M_2\!\bigl(\End^0_{\overline{\QQ}}(J(R))\bigr).
\]
对椭圆曲线，\(\End^0_{\overline{\QQ}}\) 仅可能为 \(\QQ\) 或虚二次域 \(K\)，二分即得。

再由 \(\rho(A\times B)=\rho(A)+\rho(B)\)（在 \(\Hom(A,B)=0\) 情形）、
\(
\rho(J(C))=1
\)
与椭圆平方的经典公式
\[
\rho(E^2)=
\begin{cases}
3,& E\ \text{非 CM},\\
4,& E\ \text{CM},
\end{cases}
\]
得到
\(
\rho(J(C_6)_{\overline{\QQ}})=1+\rho(J(R)^2)
\)，即 \(4/5\) 二值刚性。
\end{proof}

\endinput
